\chapter{Relative Motion}\label{ch:g10:relative-motion}

A passenger strolls toward the front of a high-speed train, coffee in
hand, at a leisurely walking pace. A trackside observer clocks that same
passenger at over \qty{300}{km/h}. Both measurements are correct. Neither
is ``the'' speed of the passenger, because there is no such thing: motion
is not a property of a body alone, but of a body \emph{and} the point of
view chosen to describe it. This chapter makes that point of view precise
--- and teaches you to pick a good one.

\section{Motion is relative}

\begin{definition}[Reference frame]\label{def:g10:relative-motion:frame}
A \emph{reference frame}\index{reference frame} is a rigid body --- the
ground, a train carriage, the Earth as a whole --- relative to which
positions are recorded, together with a clock to date them. A body is
\emph{at rest}\index{rest} in a frame when its position in that frame
does not change over time; otherwise it is in motion in that frame.
\end{definition}

\begin{definition}[Trajectory]\label{def:g10:relative-motion:trajectory}
The \emph{trajectory}\index{trajectory} of a body in a given \omterm{def:g10:relative-motion:frame}{reference
frame} is the curve formed by its successive positions in that frame: a
straight line, a circle, an arc\dots{} The trajectory only makes sense
once the frame is named.
\end{definition}

\begin{example}[The walker on the train]\label{ex:g10:relative-motion:walker}
A train travels at \qty{300}{km/h} relative to the ground while a
passenger walks toward the front at \qty{4}{km/h} relative to the train.
\begin{itemize}
  \item In the \emph{train} frame: the passenger moves at \qty{4}{km/h}
        down the aisle; the seats are at \omterm{def:g10:relative-motion:frame}{rest}; the trees outside rush
        backward at \qty{300}{km/h}.
  \item In the \emph{ground} frame: the seats move at \qty{300}{km/h},
        the passenger at \qty{304}{km/h}, and the trees are at \omterm{def:g10:relative-motion:frame}{rest}.
\end{itemize}
The same scene, two frames, two entirely different descriptions --- and
no experiment performed inside the train can declare one of them ``the
real one''.
\end{example}

\begin{proposition}[Relativity of motion]\label{prop:g10:relative-motion:relativity}
The state of \omterm{def:g10:relative-motion:frame}{rest} or motion of a body, its \omterm{def:g10:relative-motion:trajectory}{trajectory}, and its speed all
depend on the \omterm{def:g10:relative-motion:frame}{reference frame} chosen to describe them.
\end{proposition}

\begin{proof}
One example settles each claim. The seated passenger of
\cref{ex:g10:relative-motion:walker} is at \omterm{def:g10:relative-motion:frame}{rest} in the train frame and
moves at \qty{300}{km/h} in the \omterm{def:g10:relative-motion:frames}{ground frame} (\omterm{def:g10:relative-motion:frame}{rest} and speed depend on
the frame). For the \omterm{def:g10:relative-motion:trajectory}{trajectory}, watch the valve on a bicycle wheel: an
observer riding the bicycle sees it trace a circle around the hub, while
an observer standing on the road sees it trace a sequence of arches (a
curve called a \emph{cycloid}\index{cycloid}) --- see the figure below.
\end{proof}

\begin{omfigure}
\begin{tikzpicture}[scale=0.85, font=\small]
  \draw[gray!60] (-1.6,0) -- (1.6,0);
  \draw[gray, thin] (0,1) circle (1);
  \fill (0,1) circle (1.5pt);
  \draw[gray] (0,1) -- (0.71,1.71);
  \draw[omDef, thick, -Stealth] (1,1) arc (0:330:1);
  \fill[omThm] (0.71,1.71) circle (2pt) node[above right] {valve};
  \node[omDef] at (0,-0.5) {bicycle frame: a circle};
\end{tikzpicture}
\qquad
\begin{tikzpicture}[scale=0.42, font=\small]
  \draw[gray!60] (-0.6,0) -- (13.2,0);
  \draw[omDef, thick] plot[domain=0:12.566, samples=100, smooth]
    ({\x - sin(\x r)}, {1 - cos(\x r)});
  \draw[gray, thin] (9.42,1) circle (1);
  \fill (9.42,1) circle (1.5pt);
  \draw[gray] (9.42,1) -- (9.42,2);
  \fill[omThm] (9.42,2) circle (3pt);
  \node[omDef] at (6.3,-1.0) {road frame: a cycloid};
\end{tikzpicture}

{\small The same valve, two frames: a circle for the cyclist, a \omterm{prop:g10:relative-motion:relativity}{cycloid}
--- a chain of arches --- for the observer on the road.}
\end{omfigure}

\begin{remark}
Note what does \emph{not} change in the two pictures: the valve, the
wheel, the physics. Only the description changes. Asking ``is the valve
really going in a circle?'' is like asking whether a mountain is
``really'' to the left: it depends on where you stand.
\end{remark}

\section{Average speed}

\begin{definition}[Average speed]\label{def:g10:relative-motion:speed}
The \emph{average speed}\index{speed!average} of a body, in a given
\omterm{def:g10:relative-motion:frame}{reference frame}, is the distance $d$ it travels in that frame divided by
the travel time $t$:
\[
v = \frac{d}{t} .
\]
With $d$ in \omterm{def:g10:orders-of-magnitude:unit}{metres} and $t$ in seconds, $v$ is in \omterm{def:g10:orders-of-magnitude:unit}{metres} per second
(\unit{m/s}); everyday life prefers kilometres per hour (\unit{km/h}).
\end{definition}

\begin{method}[Converting between m/s and km/h]\label{met:g10:relative-motion:convert}
One kilometre is \qty{1000}{m} and one hour is \qty{3600}{s}, so
$\qty{1}{m/s} = \qty{3.6}{km/h}$:
\begin{itemize}
  \item from \unit{m/s} to \unit{km/h}: \emph{multiply} by $3.6$;
  \item from \unit{km/h} to \unit{m/s}: \emph{divide} by $3.6$.
\end{itemize}
Sanity check: the number of \unit{km/h} is always the \emph{larger} of
the two.
\end{method}

\begin{example}[Some speeds worth knowing]\label{ex:g10:relative-motion:convert}
A brisk walk: $\qty{1.5}{m/s} \times 3.6 = \qty{5.4}{km/h}$. A city car
at \qty{54}{km/h}: $54 / 3.6 = \qty{15}{m/s}$ --- about \qty{15}{m}
covered every second, worth remembering before stepping off a kerb.
Sound in air: $\qty{340}{m/s} \times 3.6 = \qty{1224}{km/h}$.
\end{example}

\begin{example}[A hike with a break]\label{ex:g10:relative-motion:hike}
A hiker covers \qty{12}{km} in \qty{2}{h} of walking, then \omterm{def:g10:relative-motion:frame}{rests} for
\qty{1}{h}, then covers \qty{6}{km} in the final \qty{1.5}{h}. \omterm{def:g10:relative-motion:speed}{Average
speed} over the whole outing:
\[
v = \frac{d}{t} = \frac{12 + 6}{2 + 1 + 1.5}
= \frac{\qty{18}{km}}{\qty{4.5}{h}} = \qty{4.0}{km/h}.
\]
The average counts \emph{all} the elapsed time, breaks included --- it is
generally not the average of the speeds of the separate legs.
\end{example}

\section{Choosing a frame: ground, geocentric, heliocentric}

Since every frame is legitimate, we choose the one in which the motion we
care about looks simplest. Three choices cover most of the physics in
this book.

\begin{definition}[The three standard frames]\label{def:g10:relative-motion:frames}
\begin{itemize}
  \item The \emph{ground frame}\index{ground frame} (or
        \emph{laboratory frame}\index{laboratory frame}) is attached to
        the Earth's surface where the experiment happens. It is the
        right frame for everyday motion: falling objects, vehicles,
        sports, laboratory benches.
  \item The \emph{geocentric frame}\index{geocentric frame} is centred
        on the Earth but does \emph{not} spin with it (its directions
        point to distant stars). It is the right frame for the Moon and
        for artificial \omterm{def:g10:universal-gravitation:satellite}{satellites}, which \omterm{def:g10:universal-gravitation:satellite}{orbit} the Earth as a whole and
        ignore its daily rotation.
  \item The \emph{heliocentric frame}\index{heliocentric frame} is
        centred on the Sun, with directions again fixed on distant
        stars. It is the right frame for the planets: each then follows
        a simple, nearly circular \omterm{def:g10:universal-gravitation:satellite}{orbit}.
\end{itemize}
\end{definition}

\begin{omfigure}
\begin{tikzpicture}[scale=1.0, font=\small]
  \fill[omProp] (0,0) circle (0.26);
  \node[below=3pt] at (0,-0.26) {Sun};
  \draw[dashed, gray] (0,0) circle (2.5);
  \draw[-Stealth, gray] (2.5,0) arc (0:22:2.5);
  \fill[omDef] (2.5,0) circle (0.12);
  \node[below right] at (2.55,-0.10) {Earth};
  \draw[dashed, gray] (2.5,0) circle (0.75);
  \draw[-Stealth, gray] (3.25,0) arc (0:55:0.75);
  \fill[gray] (2.5,0.75) circle (0.06);
  \node[above] at (2.5,0.80) {Moon};
  \draw[omThm] (2.38,0.12) -- (1.2,1.35);
  \node[omThm, left, align=right] at (1.25,1.45)
    {ground frame:\\a point spinning\\with the surface};
  \node[omDef, align=left, anchor=west] at (3.6,0.1)
    {geocentric frame:\\centred on the Earth\\(Moon, satellites)};
  \node[omProp, align=left, anchor=west] at (-1.4,-1.6)
    {heliocentric frame: centred on the Sun (planets)};
\end{tikzpicture}

{\small Three nested frames: the ground spins with the Earth, the
\omterm{def:g10:relative-motion:frames}{geocentric frame} rides with the Earth without spinning, the \omterm{def:g10:relative-motion:frames}{heliocentric
frame} watches the whole dance from the Sun.}
\end{omfigure}

\begin{method}[Choosing the frame]\label{met:g10:relative-motion:choose}
Name the bodies whose motion you want to describe, then pick the frame in
which that motion is simplest:
\begin{enumerate}
  \item motion near the Earth's surface, lasting minutes --- \omterm{def:g10:relative-motion:frames}{ground
        frame};
  \item \omterm{def:g10:universal-gravitation:satellite}{orbits} \emph{around} the Earth (Moon, \omterm{def:g10:universal-gravitation:satellite}{satellites}) --- \omterm{def:g10:relative-motion:frames}{geocentric
        frame};
  \item \omterm{def:g10:universal-gravitation:satellite}{orbits} around the Sun (planets, comets) --- \omterm{def:g10:relative-motion:frames}{heliocentric frame}.
\end{enumerate}
Always state your choice before describing a motion: a \omterm{def:g10:relative-motion:trajectory}{trajectory} or a
speed quoted without its frame is meaningless
(\cref{prop:g10:relative-motion:relativity}).
\end{method}

\begin{remark}[Who is right, Ptolemy or Copernicus?]\label{rem:g10:relative-motion:history}
``The Sun rises in the east'' and ``the Earth spins toward the east'' are
the same observation reported in two frames --- the first in the \omterm{def:g10:relative-motion:frames}{ground
frame}, where the Sun sweeps across the sky, the second in the \omterm{def:g10:relative-motion:frames}{geocentric
frame}, where the observer is carried around the Earth's axis. For
centuries the frames themselves were the battlefield: Ptolemy's
Earth-centred description tracked the sky accurately but needed loops
upon loops to follow the planets, while the Sun-centred description
revived by Copernicus in 1543 made each planet a simple near-circle and
explained the strange backward loops of Mars as a mere frame effect ---
the Earth overtaking a slower planet. The modern verdict is not that one
frame is true and the other false, but that they are not equally
\emph{convenient}; a deeper criterion for preferring some frames
(the inertial ones) arrives with the principle of inertia in
\cref{ch:g10:inertia}.
\end{remark}

\section{Composing motions along a line}

What connects the \qty{4}{km/h} of the walker (train frame) to the
\qty{304}{km/h} (\omterm{def:g10:relative-motion:frames}{ground frame})? Along a single line, the answer is plain
addition.

\begin{proposition}[Composition of speeds in one dimension]\label{prop:g10:relative-motion:compose}
Suppose a body moves along a line on a moving support (a walkway, a
train, flowing water), the support itself sliding along the same line
relative to the ground. Choosing a positive direction on the line and
counting speeds with a sign,
\[
v_{\text{body/ground}} \;=\; v_{\text{body/support}} \;+\;
v_{\text{support/ground}} .
\]
In words: speeds in the same direction add; speeds in opposite
directions subtract.
\end{proposition}

\begin{proof}
During one second, the support carries the body $v_{\text{support/ground}}$
\omterm{def:g10:orders-of-magnitude:unit}{metres} along the line, and the body advances $v_{\text{body/support}}$
\omterm{def:g10:orders-of-magnitude:unit}{metres} along the support. Both displacements happen along the same line,
so the body's total displacement over the ground in that second is their
algebraic sum --- which is its speed relative to the ground.
\end{proof}

\begin{example}[The moving walkway]\label{ex:g10:relative-motion:walkway}
An airport walkway rolls at \qty{1.0}{m/s}. A traveller walks on it at
\qty{1.5}{m/s} relative to the belt, in the direction of travel: speed
relative to the ground $1.5 + 1.0 = \qty{2.5}{m/s}$. Walking the wrong
way at the same pace: $1.5 - 1.0 = \qty{0.5}{m/s}$, still (slowly)
progressing against the belt. Standing still on the belt:
\qty{1.0}{m/s} relative to the ground, \qty{0}{m/s} relative to the
belt --- at \omterm{def:g10:relative-motion:frame}{rest} and in motion at the same time, in two different
frames.
\end{example}

\begin{omfigure}
\begin{tikzpicture}[font=\small]
  \draw[fill=gray!12] (0,0) rectangle (7,0.3);
  \node[gray] at (6.3,0.15) {belt};
  \draw (2.5,0.3) -- (2.5,1.05);
  \draw (2.5,0.75) -- (2.2,0.45) (2.5,0.75) -- (2.8,0.45);
  \fill (2.5,1.2) circle (2.5pt);
  \draw[-Stealth, omDef, thick] (0.4,-0.45) -- (1.9,-0.45)
    node[midway, below] {belt/ground: \qty{1.0}{m/s}};
  \draw[-Stealth, omThm, thick] (2.5,1.6) -- (4.75,1.6)
    node[midway, above] {walker/belt: \qty{1.5}{m/s}};
  \draw[-Stealth, omMeth, very thick] (2.5,2.3) -- (6.25,2.3)
    node[midway, above] {walker/ground: \qty{2.5}{m/s}};
\end{tikzpicture}

{\small One-dimensional composition: the arrows are drawn to scale, and
the ground-frame arrow is the sum of the other two.}
\end{omfigure}

\begin{remark}[Two directions at once]\label{rem:g10:relative-motion:twodim}
When the two motions are \emph{not} along the same line --- a boat aiming
straight across a river while the current pushes it downstream --- each
motion runs its course independently, and the velocities compose as
arrows placed tip to tail, exactly like the vectors of the mathematics
volume. The boat still reaches the far bank in the time set by its
crossing speed alone, but it lands downstream of its aiming point, and
its ground-frame \omterm{def:g10:relative-motion:trajectory}{trajectory} is a slanted straight line. The full vector
treatment comes in a later chapter; here we only read the picture.
\end{remark}

\begin{omfigure}
\begin{tikzpicture}[scale=0.95, font=\small]
  \fill[omDef!10] (0,0) rectangle (6.4,2.4);
  \draw (0,0) -- (6.4,0) (0,2.4) -- (6.4,2.4);
  \draw[-Stealth, gray] (4.8,0.35) -- (5.6,0.35);
  \node[gray] at (5.2,0.62) {current};
  \draw[dashed, omMeth] (1.5,0) -- (3.3,2.4);
  \draw[-Stealth, omThm, thick] (1.5,0.6) -- (1.5,2.0)
    node[left, pos=0.7] {boat/water};
  \draw[-Stealth, omDef, thick] (1.5,0.6) -- (2.55,0.6)
    node[below, pos=0.8] {water/ground};
  \draw[-Stealth, omMeth, very thick] (1.5,0.6) -- (2.55,2.0)
    node[right, pos=0.75] {boat/ground};
  \fill (1.5,0.6) circle (1.8pt);
\end{tikzpicture}

{\small A boat heading straight across a river: in the water frame it
goes straight across; in the \omterm{def:g10:relative-motion:frames}{ground frame} it follows the slanted dashed
line, drifting downstream as it crosses.}
\end{omfigure}

\section{Exercises}

\begin{exercise}[$\star$]\label{exo:g10:relative-motion:1}
A passenger is asleep in a train travelling at \qty{100}{km/h}. State
whether the passenger is at \omterm{def:g10:relative-motion:frame}{rest} or in motion, and at what speed,
relative to: the carriage; the ground; a tree beside the track; an
oncoming train travelling at \qty{100}{km/h} the other way. Then describe
the motion of the tree in the frame of the passenger's train.
\end{exercise}

\begin{exercise}[$\star$]\label{exo:g10:relative-motion:2}
Convert: \qty{72}{km/h} and \qty{108}{km/h} into \unit{m/s};
\qty{25}{m/s} and \qty{340}{m/s} into \unit{km/h}. Check each answer
with the sanity rule of \cref{met:g10:relative-motion:convert}.
\end{exercise}

\begin{exercise}[$\star$]\label{exo:g10:relative-motion:3}
A cyclist covers \qty{36}{km} in \qty{1}{h}~\qty{30}{min}; a sprinter
covers \qty{100}{m} in \qty{10.0}{s}. Compute both \omterm{def:g10:relative-motion:speed}{average speeds} in
\unit{m/s} and in \unit{km/h}. Who is faster, and does the answer
surprise you?
\end{exercise}

\begin{exercise}[$\star$]\label{exo:g10:relative-motion:4}
A moving walkway rolls at \qty{0.9}{m/s}. A traveller walks at
\qty{1.4}{m/s} relative to the belt. Give the traveller's speed relative
to the ground when walking in the belt's direction, against it, and when
standing still on the belt.
\end{exercise}

\begin{exercise}[$\star$]\label{exo:g10:relative-motion:5}
For each motion, choose the most convenient frame among ground,
geocentric and heliocentric, and justify in one sentence: a pen dropped
inside a moving train (bonus: an even better frame is available); the
\omterm{def:g10:universal-gravitation:satellite}{orbit} of the International Space Station; one Martian year; a
\qty{100}{m} sprint.
\end{exercise}

\begin{exercise}[$\star\star$]\label{exo:g10:relative-motion:6}
A train rolls at \qty{108}{km/h}. A passenger walks toward the front at
\qty{1.5}{m/s} relative to the train.
\begin{enumerate}
  \item Convert the train's speed to \unit{m/s}, then give the
        passenger's speed relative to the ground, walking toward the
        front and toward the rear.
  \item The carriage is \qty{60}{m} long. How long does the walk to the
        front take, and how far does the passenger travel relative to
        the ground during it?
\end{enumerate}
\end{exercise}

\begin{exercise}[$\star\star$]\label{exo:g10:relative-motion:7}
A boat's speed in still water is \qty{5.0}{m/s}; a river flows at
\qty{2.0}{m/s}. The boat travels \qty{420}{m} downstream, turns around,
and comes back.
\begin{enumerate}
  \item Compute the boat's speed relative to the bank on each leg, and
        the duration of each leg.
  \item Compute the \omterm{def:g10:relative-motion:speed}{average speed} over the round trip. Why is it less
        than \qty{5.0}{m/s}, even though the current helps on one leg
        exactly as much as it hinders on the other?
\end{enumerate}
\end{exercise}

\begin{exercise}[$\star\star$]\label{exo:g10:relative-motion:8}
Two trains pass each other on parallel tracks: train A (\qty{180}{m}
long) at \qty{25}{m/s}, train B (\qty{90}{m} long) at \qty{20}{m/s}.
\begin{enumerate}
  \item They travel in opposite directions. What is the speed of B in
        the frame of A, and how long does the complete crossing last
        (from noses meeting to tails separating)?
  \item Same questions if they travel in the same direction.
\end{enumerate}
\end{exercise}

\begin{exercise}[$\star\star$]\label{exo:g10:relative-motion:9}
A bicycle rolls along a straight road. Describe the \omterm{def:g10:relative-motion:trajectory}{trajectory} of the
tyre valve in: the frame of the cyclist; the \omterm{def:g10:relative-motion:frames}{ground frame}; and the frame
of the valve itself. Describe also the \omterm{def:g10:relative-motion:trajectory}{trajectory} of the wheel's hub in
the \omterm{def:g10:relative-motion:frames}{ground frame}. No computation --- name each curve.
\end{exercise}

\begin{exercise}[$\star\star$]\label{exo:g10:relative-motion:10}
The Earth's equator is about \qty{40000}{km} long, and the Earth spins
once in \qty{24}{h}.
\begin{enumerate}
  \item Compute the speed of a point on the equator in the \omterm{def:g10:relative-motion:frames}{geocentric
        frame}, in \unit{km/h} and in \unit{m/s}.
  \item Explain in one sentence each: why we do not feel this speed, and
        why ``the Sun rises'' and ``the Earth spins'' report the same
        observation.
\end{enumerate}
\end{exercise}

\begin{exercise}[$\star\star$]\label{exo:g10:relative-motion:11}
Standing still on a moving escalator, the ride up takes \qty{60}{s}.
Walking up the same escalator switched off takes \qty{90}{s}. How long
does walking up the moving escalator take? (Work with the speeds, not
the times.)
\end{exercise}

\begin{exercise}[$\star\star\star$]\label{exo:g10:relative-motion:12}
A driver covers the same distance twice: outbound at \qty{60}{km/h},
back at \qty{90}{km/h}.
\begin{enumerate}
  \item Show that the \omterm{def:g10:relative-motion:speed}{average speed} over the round trip is
        \qty{72}{km/h} --- not \qty{75}{km/h}.
  \item Explain where the naive answer goes wrong, and state which leg
        dominates the average.
\end{enumerate}
\end{exercise}

\begin{exercise}[$\star\star\star$]\label{exo:g10:relative-motion:13}
Two cyclists start \qty{35}{km} apart and ride toward each other, one at
\qty{15}{km/h}, the other at \qty{20}{km/h}.
\begin{enumerate}
  \item Describe the motion of the second cyclist in the frame of the
        first, and deduce when they meet.
  \item How far has each ridden at the meeting?
  \item If instead they ride in the same direction (the faster one
        behind), how long does the chase last?
\end{enumerate}
\end{exercise}

\begin{exercise}[$\star\star\star$]\label{exo:g10:relative-motion:14}
A swimmer crosses a river \qty{60}{m} wide, heading straight across at
\qty{1.2}{m/s} relative to the water; the current flows at
\qty{0.9}{m/s}.
\begin{enumerate}
  \item Explain why the crossing time is set by the swimmer's crossing
        speed alone, and compute it.
  \item How far downstream does the swimmer land?
  \item Using a right triangle, find the swimmer's speed in the \omterm{def:g10:relative-motion:frames}{ground
        frame} and the length of the path actually swum over the ground.
\end{enumerate}
\end{exercise}

\begin{exercise}[$\star\star\star$]\label{exo:g10:relative-motion:15}
The Earth \omterm{def:g10:universal-gravitation:satellite}{orbits} the Sun on a near-circle of radius
\qty{1.50e11}{m}, in one year ($\num{3.16e7}$ seconds).
\begin{enumerate}
  \item Compute the Earth's speed in the \omterm{def:g10:relative-motion:frames}{heliocentric frame}, in
        \unit{km/s}.
  \item Compare it with the equator's spin speed found in
        \cref{exo:g10:relative-motion:10}.
  \item Seen from the Earth, Mars occasionally seems to stop and drift
        backward among the stars for a few weeks. Explain in two
        sentences, using this chapter's vocabulary, why the \omterm{def:g10:relative-motion:frames}{heliocentric
        frame} dissolves the mystery.
\end{enumerate}
\end{exercise}

\section{Problem: The walkway, the river and the satellite}

\begin{problem}[{Weekend problem --- the moving walkway, the river and
the satellite: who moves, relative to what, and the armchair traveller's
final speed report}]
\label{pb:g10:relative-motion:1}
An airport, a ferry crossing, and a bright dot sliding across the dusk
sky: three everyday scenes, one question asked three times --- who
moves, relative to what? Each part picks its frame, computes, and
watches the answer change with the point of view; the finale totals up
how fast \emph{you} are moving right now, sitting perfectly still.

\medskip
\noindent\textbf{Part I --- The walkway race.}
An airport walkway is \qty{90}{m} long and rolls at \qty{1.0}{m/s}. A
traveller walks at \qty{1.5}{m/s} (relative to whatever is underfoot).
\begin{enumerate}
  \item How long does the walkway take to carry a standing traveller
        from end to end?
  \item How long does the traveller take when walking on the moving
        belt, in its direction?
  \item How long when walking beside the walkway, on the fixed floor?
  \item A prankster walks the wrong way on the belt, at \qty{1.5}{m/s}
        relative to it. Speed relative to the ground, and time to get
        from the far end back to the entrance?
  \item A child runs the wrong way at exactly \qty{1.0}{m/s} relative
        to the belt. Describe the child's motion in the \omterm{def:g10:relative-motion:frames}{ground frame},
        and name the gym machine this reinvents.
  \item Summarize questions 1--5 in a two-column table: each mover's
        speed in the belt frame and in the \omterm{def:g10:relative-motion:frames}{ground frame}. Which column
        do the travellers' legs feel?
\end{enumerate}

\medskip
\noindent\textbf{Part II --- The river.}
A ferry crosses a river \qty{120}{m} wide. Its speed relative to the
water is \qty{2.0}{m/s}; the current flows at \qty{1.5}{m/s}. The pilot
first aims straight at the opposite bank.
\begin{enumerate}[resume]
  \item How long does the crossing take?
  \item How far downstream does the ferry land?
  \item Using a right triangle, compute the ferry's speed in the \omterm{def:g10:relative-motion:frames}{ground
        frame} and the length of its ground-frame \omterm{def:g10:relative-motion:trajectory}{trajectory}. (A
        familiar triple of numbers appears.)
  \item Describe the ferry's \omterm{def:g10:relative-motion:trajectory}{trajectory} in the water frame and in the
        \omterm{def:g10:relative-motion:frames}{ground frame}. What is the same about them, and what differs?
  \item A swimmer whose speed relative to the water is \qty{1.0}{m/s}
        sets off from the same bank, aiming straight across like the
        ferry. How long does the crossing take, and how far downstream
        does the swimmer land? Compute the swimmer's ground-frame
        speed (right triangle again). Why does the slower swimmer
        drift farther than the ferry?
  \item A log floats past the ferry. Give its speed in the \omterm{def:g10:relative-motion:frames}{ground frame}
        and in the water frame. In which frame is the log ``moving''?
        Comment.
\end{enumerate}

\medskip
\noindent\textbf{Part III --- The \omterm{def:g10:universal-gravitation:satellite}{satellite} pass.}
At dusk, the International Space Station crosses the sky in a few
minutes. It \omterm{def:g10:universal-gravitation:satellite}{orbits} at an altitude of \qty{400}{km}; the Earth's radius
is \qty{6.37e6}{m}, and one \omterm{def:g10:universal-gravitation:satellite}{orbit} takes \qty{5.56e3}{s} (about
\qty{93}{min}).
\begin{enumerate}[resume]
  \item In which standard frame is the station's motion simplest?
        Compute its speed in that frame, in \unit{km/s}.
  \item Explain in two sentences why the \omterm{def:g10:relative-motion:frames}{ground frame} describes the
        station awkwardly --- what happens to its track over the ground
        between one \omterm{def:g10:universal-gravitation:satellite}{orbit} and the next?
  \item Compare the station's orbital speed with the equator's spin
        speed ($\approx \qty{0.46}{km/s}$): by what factor is the
        station faster?
  \item A geostationary \omterm{def:g10:universal-gravitation:satellite}{satellite} \omterm{def:g10:universal-gravitation:satellite}{orbits} at radius \qty{4.22e7}{m} in
        exactly one day ($\qty{8.64e4}{s}$). Compute its speed in the
        \omterm{def:g10:relative-motion:frames}{geocentric frame}, then describe its motion in the \omterm{def:g10:relative-motion:frames}{ground frame}.
        Reconcile the two answers in one sentence.
  \item In the \omterm{def:g10:relative-motion:frames}{geocentric frame} the Moon traces a near-circle around
        the Earth in about $27$ days. Sketch in words its \omterm{def:g10:relative-motion:trajectory}{trajectory} in
        the \omterm{def:g10:relative-motion:frames}{heliocentric frame} over a few months.
\end{enumerate}

\medskip
\noindent\textbf{Part IV --- Who moves?}
\begin{enumerate}[resume]
  \item ``The Sun rises in the east.'' ``The Earth spins toward the
        east.'' State the frame in which each sentence describes the
        observation, and explain why neither is wrong.
  \item In one sober paragraph: what did Ptolemy's frame do well, what
        did Copernicus's frame simplify (mention Mars's backward
        loops), and why is the modern verdict about \emph{convenience}
        rather than about one frame being true?
  \item The armchair traveller's speed report. You sit still in an
        armchair near the equator. Give your speed in the \omterm{def:g10:relative-motion:frames}{ground frame},
        in the \omterm{def:g10:relative-motion:frames}{geocentric frame} ($\approx \qty{0.46}{km/s}$), and in
        the \omterm{def:g10:relative-motion:frames}{heliocentric frame} ($\approx \qty{30}{km/s}$). Then compute
        the distance you cover in the \omterm{def:g10:relative-motion:frames}{heliocentric frame} during a
        \qty{90}{min} film, and conclude the problem with the chapter's
        one-sentence moral.
\end{enumerate}
\end{problem}
